\section{The measurement problem}
\label{sec:problem}

Suppose a benchmark reports that a detector identifies deceptive model outputs with 95\% accuracy. To produce that number, someone had to \emph{make} a model behave deceptively, typically by instructing it to. The benchmark therefore compares two populations of text that differ in at least two ways: the presence of deception, and the presence of an instruction that induces it. A detector can score 95\% by measuring either one.

This is not specific to deception. Any behavioral benchmark that elicits a property by prompting for it inherits the same structure: sycophancy elicited by an agreeable-user prompt, sandbagging elicited by an underperform instruction, persuasion elicited by a persuade instruction. In each case the elicitation is a treatment applied to one arm and not the other, so every downstream consequence of it (register, hedging, verbosity, refusal language) is perfectly correlated with the label. The reported score therefore conflates construct sensitivity with sensitivity to the elicitation condition, and cannot be read as a clean estimate of construct measurement.

The problem is an identification problem, and it has a standard shape. Writing $C$ for the construct (was the model deceptive?), $E$ for the elicitation condition (was it instructed to be?), and $\hat{y}$ for detector output, a benchmark that varies $C$ and $E$ together cannot separate $\Pr(\hat{y}\mid C)$ from $\Pr(\hat{y}\mid E)$. The headline number identifies their union (Figure~\ref{fig:protocol}).

\paragraph{What this paper establishes.} It is tempting to read a large elicitation effect as proof that a detector exploits a shortcut. That inference does not go through, and we do not make it: the elicitation \emph{is} the mechanism inducing the target behavior, so $E$ is simultaneously the treatment creating deception and a cue correlated with the label. A detector responding to $E$ may be responding to deception, to behavioral consequences of the instruction, or to both, and \textbf{the design cannot tell them apart}. That is the failure we demonstrate: an \emph{identification} failure, not a demonstrated detector failure. The benchmark cannot license a construct-level reading of its own score; we do not claim no deception-specific signal exists.

\paragraph{Contribution.} Three levels, in order of generality. \textbf{(1) Conceptual:} behavioral benchmarks that manufacture the behavior they measure have a built-in elicitation/construct identification problem. \textbf{(2) Methodological:} three diagnostics that test whether a reported score survives it (two causal-validity checks and one generalization check), runnable on the benchmark's own claim set with no new model training (the reliability study in §\ref{sec:validity} does need human annotation). \textbf{(3) Empirical:} a deception benchmark that fails the two causal checks and passes the generalization one. We also report two stress tests (§\ref{sec:stability}) whose results a single-model evaluation would have gotten wrong.

Three claims are in play and we keep them apart throughout. \textbf{A}, that standard instructed protocols confound the construct with the elicitation so their scores do not identify the construct: strongly supported. \textbf{B}, that the detectors we audit are substantially responsive to elicitation-associated behavior: strongly supported. \textbf{C}, that no deception-specific signal exists: \emph{we do not claim this and our design cannot establish it}, since the benchmark never operationalizes deception independently of the instruction.

\begin{figure}[t]
\centering
\begin{tikzpicture}[
  box/.style={draw, rounded corners=2pt, align=center, inner sep=2.5pt, font=\footnotesize},
  lbl/.style={font=\scriptsize, align=center},
  ar/.style={-{Latex[length=3.5pt]}, semithick},
]
% --- confounded benchmark ---
\node[lbl, anchor=west, font=\scriptsize\bfseries] at (-0.15,1.20) {Standard benchmark};
\node[box, fill=black!5] (content) at (1.35,0.72) {content $V$};
\node[box, fill=black!5] (elicit)  at (1.35,0.28) {elicitation $E$};
\node[box] (model) at (3.5,0.5)  {model output};
\node[box] (det)   at (5.5,0.5)  {detector};
\node[box] (score) at (7.15,0.5) {score};
\draw[ar] (content.east) -- (model.west);
\draw[ar] (elicit.east)  -- (model.west);
\draw[ar] (model) -- (det);
\draw[ar] (det) -- (score);
\node[lbl, anchor=west, text width=2.9cm] at (7.8,0.5)
  {$V,E$ vary together:\\score identifies\\their \emph{union}};

% --- audit ---
\node[lbl, anchor=west, font=\scriptsize\bfseries] at (-0.15,-0.55) {Audit protocol};
\node[box, fill=blue!6]  (d1) at (1.9,-1.2) {\textbf{1.} equalize $E$};
\node[box, fill=blue!6]  (d2) at (4.6,-1.2) {\textbf{2.} factorial $V\!\times\!E$};
\node[box, fill=green!8] (d3) at (7.6,-1.2) {\textbf{3.} hold out groups};
\draw[ar] (d1) -- (d2);
\draw[ar] (d2) -- (d3);
\node[lbl, text width=2.6cm] at (1.9,-1.78) {does the score survive?};
\node[lbl, text width=2.9cm] at (4.6,-1.78) {$\beta_E$ vs.\ $\beta_V$: which does it track?};
\node[lbl, text width=2.9cm] at (7.6,-1.78) {does it generalize to unseen claims?};
\draw[decorate, decoration={brace, amplitude=3pt, mirror}, black!55]
  (0.75,-2.12) -- (5.95,-2.12);
\node[lbl, anchor=north] at (3.35,-2.24) {causal validity};
\draw[decorate, decoration={brace, amplitude=3pt, mirror}, black!55]
  (6.25,-2.12) -- (8.95,-2.12);
\node[lbl, anchor=north] at (7.6,-2.24) {generalization};
\end{tikzpicture}
\caption{A behavioral benchmark varies content and elicitation together, so its score identifies their union rather than the construct it names. The diagnostics cross the two factors and test what survives; they do not recover a construct estimate the design never identified.}
\label{fig:protocol}
\end{figure}


\paragraph{Case study and scope.} Our worked example is behavioral deception detection in instructed roleplay, evaluated on seven open-weight instruction-tuned LLMs from 3B to 70B plus a frontier API model, with 50 matched true/false claim pairs. We audit two mechanistically unrelated detectors: a five-feature LLM-scored behavioral pipeline, and a from-scratch reimplementation of the published mechanism of \citet{pacchiardi2023catch} (fixed unrelated yes/no elicitation probes $\to$ binary vector $\to$ logistic regression). Everything is English-language instructed roleplay; we make no claim about autonomous strategic deception \citep{hubinger2019risks, scheurer2024strategic, greenblatt2024alignment}.

\paragraph{Related work.} Behavioral lie detection via elicitation questions was introduced by \citet{pacchiardi2023catch}; white-box variants probe internal representations \citep{azaria2023internal, burns2022discovering, zou2023representation, li2024inference}. Instruction-following shapes LLM outputs in ways that confound downstream measurement \citep{sharma2023sycophancy, denison2024sycophancy, perez2022red}, and models behave differently when they infer they are being evaluated \citep{hubinger2024sleeper, greenblatt2024alignment}.

The gap between an observed score and the latent property it is taken to measure is the classical problem of \emph{construct validity} \citep{cronbach1955construct}, imported into ML measurement by \citet{jacobs2021measurement}; behavioral test suites \citep{ribeiro2020beyond} and underspecification analyses \citep{damour2020underspecification} attack neighbouring versions of it. Our identification argument is likewise standard: a factorial design separating a treatment from a co-varying condition is textbook. \textbf{The contribution is not the machinery but its absence in practice.} Behavioral LLM benchmarks routinely instantiate this exact confound (they must manufacture the behavior they measure, and the manufacturing is detectable) yet report scores without any control that would separate the two. We show the confound is material rather than theoretical (it reverses the headline conclusion on a real paradigm), that the controls are runnable on the benchmark's own claim set without new annotation, and that whether they fire is not predictable in advance: here two fire and one passes.
